\section{Билет 72. Уравнение Ван-дер-Ваальса для реальных газов. Анализ изотерм Ван-дер-Ваальса. Изотермы реального газа. Перегретая жидкость, пересыщенный пар. Область равновесных состояний вещества на диаграмме pV.}

\begin{definition}
\textbf{Уравнение Ван-дер-Ваальса} для одного моля газа:
\begin{equation}
    \left(p + \frac{a}{V^2}\right)(V - b) = R T. \label{eq:vdw_72}
\end{equation}
Анализ его изотерм показывает, что при $T < T_{\text{кр}}$ кривая имеет волнообразный вид. Физически реализуемый процесс фазового перехода (испарение/конденсация) происходит при постоянном давлении $p_{\text{н}}$ (давление насыщенного пара).
\end{definition}

\begin{theorem}[Правило Максвелла]
Для замены неустойчивого участка изотермы Ван-дер-Ваальса на горизонтальную площадку фазового равновесия применяется \textbf{правило равных площадей Максвелла}:
горизонтальная прямая $p = p_{\text{н}}$ проводится так, чтобы площади криволинейных треугольников, отсекаемых ею от S-образной петли над и под прямой, были равны:
\begin{equation}
    \int_{V_{\text{ж}}}^{V_{\text{г}}} (p(V) - p_{\text{н}}) \, dV = 0. \label{eq:maxwell_rule_72}
\end{equation}
Это условие следует из требования равенства химических потенциалов жидкости и пара в равновесии и гарантирует, что работа, совершаемая при циклическом процессе вдоль петли, равна нулю.
\end{theorem}

\begin{definition}
\textbf{Изотермы реального газа} в координатах $(p, V)$ имеют характерный вид: при $T > T_{\text{кр}}$ монотонно убывают; при $T < T_{\text{кр}}$ содержат горизонтальный участок, соответствующий сосуществованию жидкости и пара при постоянном давлении насыщенного пара. Концы горизонтального участка соответствуют объёмам насыщенной жидкости $V_{\text{ж}}$ и насыщенного пара $V_{\text{г}}$.
\end{definition}

\begin{definition}
\textbf{Метастабильные состояния} --- состояния термодинамического неравновесия, существующие относительно длительное время при отсутствии центров фазового перехода:
\begin{itemize}
    \item \textbf{Перегретая жидкость} --- жидкость, нагретая выше температуры кипения при данном давлении, но не кипящая из-за отсутствия центров парообразования (пылинок, шероховатостей).
    \item \textbf{Пересыщенный пар} --- пар, охлаждённый ниже температуры конденсации, но не конденсирующийся из-за отсутствия центров конденсации (ионов, пылинок).
\end{itemize}
Эти состояния соответствуют участкам изотермы Ван-дер-Ваальса между точкой насыщения и экстремумами (максимумом/минимумом). Они устойчивы к малым возмущениям, но переход в стабильное состояние происходит скачком при появлении зародышей новой фазы.
\end{definition}

\begin{remark}
\textbf{Область равновесных состояний на диаграмме $(p, V)$.}
Под критической изотермой на $pV$-диаграмме существует область, ограниченная \textbf{бинодалью} (кривой насыщения). Внутри этой бинодали вещество существует в виде двухфазной системы (жидкость + пар).
\begin{itemize}
    \item Левая ветвь бинодали --- кривая насыщенной жидкости.
    \item Правая ветвь бинодали --- кривая насыщенного пара.
    \item В точке пересечения ветвей находится \textbf{критическая точка} $(p_{\text{кр}}, V_{\text{кр}}, T_{\text{кр}})$.
\end{itemize}
При пересечении бинодали изотерма идёт горизонтально. За пределами бинодали (при высоких $T$ или низких $p/V$) вещество находится в однофазном состоянии (газ или жидкость).
\end{remark}
